\begin{proposition}[Sections on a multiple abelian fibre]\hypertarget{tag-01AM}{}\label{cdp:st:A.3.2}
\StatementRecord{cdp:st:A.3.2}{proved}{01AM}{the definitions and constructions stated in the proof}

There is a compact smooth complex threefold $g:Y\to\PP^1$ with connected general fibre an abelian surface and a smooth fibre support $S$ of multiplicity $m$, for every $m\ge3$, and a line bundle $B$ on $Y$, such that
\[
H^0(S,(\Omega_Y^1\otimes B)|_S)=0,
\qquad h^0(mS,(\Omega_Y^1\otimes B)|_{mS})=3.
\tag{A.12}\hypertarget{tag-0041}{}\label{cdp:eq:A.12}
\]
Its conormal has exact order $m$.
\end{proposition}


\begin{proof}

\begin{enumerate}[label=\arabic*.,leftmargin=*]
\item Take an abelian surface $F$ and a translation $a:F\to F$ of exact order $m$. Let $\zeta$ be a primitive $m$th root of unity. The action generated by $(p,z)\mapsto(a(p),\zeta z)$ on $F\times\PP^1$ is free, so the quotient $Y$ is a smooth compact threefold and the quotient map is finite \'etale. The invariant map $t=z^m$ gives $g:Y\to\PP^1$. Away from $0$ and $\infty$ its fibre is isomorphic to $F$: the $m$ lifts of a base point are permuted transitively and freely. At $0$ the fibre is $mS$, where $S=F/\langle a\rangle$. Set $B=\OO_Y(-S)$ and $E=\Omega_Y^1\otimes B$.

\item The connected cyclic cover $F\to S$ identifies $N^*=B|_S$ with the character line of weight one. Its order is exactly $m$: otherwise a nontrivial proper character power would have a nonzero invariant trivialization on the connected compact cover, whose holomorphic functions are constant. Pullback to $F$ splits the ambient cotangent bundle into two invariant tangent forms and the normal form $dz$ of weight one. Consequently
\[
 (\Omega_Y^1\otimes B)|_S\simeq (N^*)^{\oplus2}\oplus(N^*)^2.
\]
The nontrivial flat lines on the abelian surface $S$ have no holomorphic sections. Their pullbacks to $F$ are trivial; a section would be a constant function transforming by a nontrivial character, which is zero. Since $m\ge3$, both characters in the preceding splitting are nontrivial, proving the first equality.

\item The inverse image of $mS$ is the ringed space with structure sheaf $\OO_F\otimes\C[z]/(z^m)$. Write $e$ for the local generator of the pulled-back line $B$; it has weight one. All global sections on this inverse image are linear combinations of $z^k\alpha_1e,z^k\alpha_2e,z^kdz\,e$ for $0\le k<m$, where $\alpha_1,\alpha_2$ are a basis of invariant one-forms on $F$. The first two have weight $k+1$ and the third weight $k+2$. The invariant monomials are exactly
$z^{m-1}\alpha_1e,z^{m-1}\alpha_2e,z^{m-2}dz\,e$.
Taking invariants is exact over $\C$, by averaging over the finite group. These three monomials form a basis. Translations act identically on $H^q(F,\OO_F)$, whose dimensions are $\binom{2}{q}$. The same calculation gives all fibre groups: $h^q(S,E|_S)=(0,0,0)$ and $h^q(mS,E|_{mS})=(3,6,3)$ for $q=0,1,2$. On the compact threefold the two tangent summands pull back to $\OO_F\boxtimes\OO_{\PP^1}(-1)$ and have no cohomology. The normal summand is $\OO_F\boxtimes\Omega^1_{\PP^1}(-\{0\})$. Its first cohomology on $\PP^1$ has the \v Cech basis $dz/z,dz$, of weights zero and one, respectively. Indeed sections near zero are multiples of $z\,dz$, and sections near infinity are multiples of $z^{-2}dz$, leaving exactly those two Laurent classes. Its invariant part is one-dimensional. K\"unneth and finite-cover descent therefore give $h^q(Y,E)=(0,1,2,1)$ for $q=0,1,2,3$.

\end{enumerate}
\end{proof}